\section{Emperor Theodosius: Callinicum and Thessalonica}

In 387, Ambrose undertook a second embassy to the usurper Maximus at Trier on behalf of Valentinian~II. This time the mission failed; Ambrose refused ecclesiastical communion with Maximus as the murderer of Emperor Gratian, whereupon Maximus expelled him from Trier and later in the year crossed the Alps into Italy (\work{Prolegomena}~\S III, under 387, noting \work{Ep.}~24). Justina and Valentinian fled east to the court of Emperor Theodosius~I, who crossed the Alps in 388, defeated and executed Maximus at Aquileia, restored Valentinian~II (now reconciled to Nicene orthodoxy after Justina's death), and took up residence in Milan for the next three years.

Theodosius~I was a baptized Nicene Catholic, a veteran soldier of genuine piety, and an emperor of explosive temper. His presence in Ambrose's cathedral city produced the two most famous and most starkly contrasting confrontations between bishop and emperor in late antiquity: Callinicum in 388, and Thessalonica in 390.

\subsection{Callinicum (388): Zeal, Power, and the Cost of Coercion}

Late in 388, while Ambrose was away at Aquileia, a report reached Theodosius from the military Count of the East concerning two acts of Christian mob violence in Mesopotamia. At Callinicum on the Euphrates, a local crowd---reportedly incited by the local bishop---had burned the Jewish synagogue; nearby, Christian monks angered when Valentinian gnostics blocked a psalm-singing procession on the feast of the Maccabees had burned the Valentinians' rural meetinghouse (\work{Letter} XL.6, 16; \work{Letter} XLI.1). Acting on Roman public-order law, Theodosius issued an entirely standard administrative order: the offending monks were to be punished, stolen synagogue funds investigated, and the bishop of Callinicum required to rebuild the synagogue at his own expense.

Ambrose responded with \work{Letter} XL, one of the most fiercely argued---and to modern Christian conscience, most deeply troubling---documents in his corpus. Claiming the prophet's duty of fearless speech before kings (\work{Letter} XL.2--4), he argued that forcing a Catholic bishop to finance a synagogue placed him in an intolerable dilemma between apostasy and martyrdom (XL.6--7). Remarkably, Ambrose offered to assume the arsonist's guilt rhetorically himself:

\begin{quote}
I declare that I set fire to the synagogue, or at least that I ordered those who did it, that there might not be a place where Christ was denied. \ldots\ Shall the Jews write this inscription on the front of their synagogue: ``The temple of impiety, erected from the plunder of Christians''? But, perhaps, the cause of discipline moves you, O Emperor. Which, then, is of greater importance, the show of discipline or the cause of religion? It is needful that judgment should yield to religion. (\work{Letter} XL.8, 10--11)
\end{quote}

Ambrose piled up every polemical argument at his disposal: he noted that when churches were burned by pagans and Jews under Emperor Julian at Damascus, Gaza, Ascalon, Berytus, and Alexandria, no imperial restitution had been exacted (XL.15); he called the synagogue ``a home of unbelief, a house of impiety, a receptacle of folly'' (XL.14); he warned that punishing Christian rioters would hand Jews a ``triumph over Christ's people'' (XL.20); and he pointed to the defeated usurper Maximus, who had issued an edict protecting a burned Roman synagogue shortly before his fall (XL.23).

When the written letter failed to secure a full cancellation, Ambrose carried the contest directly into the liturgy of Milan, as he described to Marcellina in \work{Letter} XLI. Preaching before Theodosius on Jeremiah's almond rod and the sinful woman who washed Christ's feet in Simon the Pharisee's house (Luke 7:36--50), Ambrose wove a sustained supersessionist contrast between Church and Synagogue before turning directly to the emperor in the nave: ``Wherefore, O Emperor, that I may now address my words not only about you, but to you \ldots\ honour the lowest, so that the angels may rejoice in their forgiveness'' (\work{Letter} XLI.26).

What followed at the foot of the pulpit is preserved in Ambrose's own words:

\begin{quote}
When I came down from the pulpit, he said to me: ``You spoke about me.'' I replied: ``I dealt with matters intended for your benefit.'' Then he said: ``I had indeed decided too harshly about the repairing of the synagogue by the bishop, but that has been rectified. The monks commit many crimes.'' Then Timasius the general began to be over-vehement against the monks, and I answered him: ``With the Emperor I deal as is fitting, because I know that he has the fear of God, but with you, who speak so roughly, one must deal otherwise.'' Then, after standing for some time, I said to the Emperor: ``Let me offer for you without anxiety, set my mind at ease.'' As he continued sitting and nodded, but did not give an open promise, and I remained standing, he said that he would amend the edict. I went on at once to say that he must end the whole investigation, lest the Count should use the opportunity of the investigation to do any injury to the Christians. He promised that it should be so. I said to him, ``I act on your promise,'' and repeated, ``I act on your promise.'' ``Act,'' he said, ``on my promise.'' And so I went to the altar, whither I should not have gone unless he had given me a distinct promise. (\work{Letter} XLI.27--28)
\end{quote}

Candor requires weighing Callinicum without either cynicism or special pleading. Ambrose genuinely feared that military inquiries in the East would subject ordinary Christians and bishops to brutal reprisals, and Theodosius himself conceded on the spot that compelling a bishop personally to rebuild a synagogue had been ``too harshly'' ordered. Yet Ambrose went far beyond seeking proportionate clemency for the bishop: he pressured the emperor into aborting the entire criminal investigation, subordinated Roman civic order (\latin{disciplina}) to religious triumphalism, and used harsh anti-Jewish rhetoric that excused mob destruction of a legally recognized minority's house of worship. Catholic reception honors Ambrose as saint and Doctor for his holiness, Nicene fidelity, and pastoral doctrine; neither sanctity nor the doctorate canonizes every political tactic, polemical excess, or legal judgment of a fourth-century bishop.

\subsection{Thessalonica (390): Bloodshed, Nathan's Rebuke, and Imperial Penance}

Two years later, the same fearless refusal to bend the altar to imperial will stood in defense of innocent human blood. In the spring of 390, a riot erupted in Thessalonica over the arrest of a popular charioteer; the crowd lynched the Gothic military commander Butheric and several garrison officers. Enraged, Theodosius ordered a reprisal---and though he later sent a countermand too late to stop it, troops summoned citizens into the circus and massacred thousands (reported in later tradition as roughly seven thousand, \work{Letter} LI, n.~3618).

Word of the atrocity reached Milan just as a synod of bishops, joined by visitors from Gaul, had assembled there. Every bishop lamented the slaughter; Ambrose realized that the emperor's fellowship with Milan would make the bishop complicit unless reconciliation were demanded (\work{Letter} LI.6). Theodosius, stung earlier by leaks from his consistory, had already cut Ambrose off from court intelligence (LI.2). Rather than staging a public collision that might inflame the emperor's pride, Ambrose withdrew from Milan---citing real illness as well as pastoral prudence---and wrote \work{Letter} LI entirely in his own hand, marked for Theodosius's eyes alone (LI.5, 14).

Ambrose began with remarkable psychological tenderness, acknowledging Theodosius's faith and repeated acts of clemency while naming his fatal weakness: ``you have a natural vehemence, which, if any one endeavours to soothe, you quickly turn to mercy; if any one stirs it up, you rouse it so much more that you can scarcely restrain it'' (\work{Letter} LI.4). Then he spoke as Nathan to King David:

\begin{quote}
There was that done in the city of the Thessalonians of which no similar record exists, which I was not able to prevent happening; which, indeed, I had before said would be most atrocious when I so often petitioned against it, and that which you yourself show by revoking it too late you consider to be grave, this I could not extenuate when done. \ldots\ Are you ashamed, O Emperor, to do that which the royal prophet David, the forefather of Christ, according to the flesh, did? \ldots\ You are a man, and it has come upon you, conquer it. Sin is not done away but by tears and penitence. Neither angel can do it, nor archangel. The Lord Himself, Who alone can say, ``I am with you,'' if we have sinned, does not forgive any but those who repent. (\work{Letter} LI.6--7, 11)
\end{quote}

Ambrose then stated his pastoral decision plainly, adding that a dream the night before his departure had confirmed it:

\begin{quote}
I dare not offer the sacrifice if you intend to be present. Is that which is not allowed after shedding the blood of one innocent person, allowed after shedding the blood of many? I do not think so. Lastly, I am writing with my own hand that which you alone may read. \ldots\ You shall then make your offering when you have received permission to sacrifice, when your offering shall be acceptable to God. \ldots\ ``I prefer mercy,'' it is said, ``rather than sacrifice.'' Are they not, then, rather Christians in truth who condemn their own sin, than they who think to defend it? (\work{Letter} LI.13--15)
\end{quote}

\witnesslimit{Ambrose's own contemporary \work{Letter} LI is a private, handwritten appeal urging Theodosius to abstain voluntarily from the Eucharist until he had performed penance. It contains no physical confrontation at the cathedral porch. The celebrated dramatic scene of Ambrose standing in the Milanese church doorway to bar the approaching emperor's entry comes from the fifth-century Eastern historian Theodoret of Cyrrhus (\work{Historia Ecclesiastica} V.18). The historical core---that Ambrose excluded Theodosius from Communion after Thessalonica and that Theodosius submitted to public penance before reconciliation---is firmly attested by Ambrose himself and early witnesses; the doorway tableau belongs to later historical dramatization.}

Theodosius yielded. Stripped of imperial insignia during his period of penance, he sought reconciliation with the Church and was readmitted to Communion---an act of royal humility that Ambrose later praised in his funeral oration for Theodosius (\work{De Obitu Theodosii}, summarized in de Romestin's Prolegomena). Unlike Callinicum, where Ambrose used Eucharistic leverage to shield arsonists from Roman law, at Thessalonica he placed the moral authority of the Gospel directly between an autocrat and the massacre of his own unarmed subjects, proving that the purple gave no exemption from the fifth commandment or the tears of repentance.
